\documentclass{article}
\usepackage[utf8]{inputenc}
\usepackage{xcolor}
\usepackage{hyperref}

\usepackage{rotating}

\usepackage{tabls}
\usepackage{booktabs}
\usepackage{amsmath}
\usepackage{amssymb}
\usepackage{amsthm}
\usepackage{amsfonts}
\usepackage{multicol}
\usepackage{enumitem}
\usepackage{microtype}
\usepackage{tikz}
\usepackage{pgfplots}
\usetikzlibrary{positioning}
\usepackage{multirow}

\usepackage{comment}

\usepackage{graphicx}
\usepackage{wrapfig}

\theoremstyle{plain}
\newtheorem{theorem}{Theorem}
\newtheorem*{theorem*}{Theorem}
\newtheorem{lemma}[theorem]{Lemma}
\newtheorem*{lemma*}{Lemma}
\newtheorem{proposition}[theorem]{Proposition}
\newtheorem{corollary}[theorem]{Corollary}

\theoremstyle{box}
\newtheorem{definition}[theorem]{Definition}
\newtheorem*{axiom*}{Axiom}
\newtheorem{dfn}[theorem]{Definition}
\newtheorem*{definition*}{Definition}
\newtheorem*{exercise*}{Exercise}
\newtheorem{observation}[theorem]{Observation}
\newtheorem{remark}[theorem]{Remark}
\newtheorem{example}[theorem]{Example}
\newtheorem{question}[theorem]{Question}
\newtheorem*{prep-problems}{Preparation Problems}

\newcommand{\pd}[3]{\frac{\partial^{#3} {#1}}{\partial {#2}^{#3}}}
\newcommand{\fpd}[2]{\frac{\partial {#1}}{\partial {#2}}}
\newcommand{\diff}[2]{\frac{d {#1}}{d {#2}}}


\begin{document}
\noindent One of the objectives of Math 119 is to deepen our understanding of functions.
\begin{definition*}[Function]\hfill
\begin{center}A \underline{function from \textbf{set D} to a \textbf{set Y}} is\\[0.25in]
 \textbf{a rule}\\[0.25in]
that assigns each element in \textbf{set D} exactly one element from \textbf{set Y}.\end{center}
\end{definition*}
\noindent BE AWARE: Often in algebra and calculus the two sets are clear from context and we take a short cut and treat the rule (or expression) as the function.
\vspace{0.1in}

\begin{definition*}[Domain]\hfill \\
The \underline{domain} of a function from \textbf{set D} to a set Y is the set of inputs of the function (the set D).
\end{definition*}

\begin{definition*}[Range]\hfill \\
The \underline{range} of a function is the set of outputs of the function.
\end{definition*}
\vspace{0.1in}

\newpage
\noindent There are at least four ways to represent a function:
\begin{itemize}
\item Words -- A Story. (verbal)
\item Algebraic (symbolic)
\item Graph (visual)
\item Table (numeric)
\end{itemize}
We need to develop the ability to use all these ways of representing a function and to be able to go back and forth between representations to gather information about the function.

\newpage
\begin{definition*}[Implied Domain]
When a function described by an equation (algebraically or symbolically), the \underline{implied domain} of the function is the set of all real values. Some people call the implied domain the domain of the function.
\end{definition*}

\textbf{Exercises}: Practice with these definitions.
\begin{itemize}
\item Find the implied domain of the function $f(x) = \sqrt{x-3}$.
\item Find the range of the function $f(x) = \sqrt{x-3}$.
\item Find the implied domain of the function $g(x) = \frac{3x-2}{x-5}$.
\item Find the range of the function$g(x) = \frac{3x-2}{x-5}$.
\item Find the implied domain of the function $h(x) = 3x^2 +4$.
\item Find the range of the function $h(x) = 3x^2 +4$.
\end{itemize}
%Adapted from Elaine (Function instruction and homework.)
\vspace{0.1in}

\newpage
\begin{definition*}[Describing Functions]\hfill
\begin{itemize}
\item We say a function $f$ is \underline{constant on the interval $I$} if for all $x_1, x_2$ in the interval, $f(x_1)=f(x_2)$.
\item We say a function $f$ is \underline{increasing on the interval $I$} if for all $x_1, x_2$ in the interval, $f(x_1) \leq f(x_2)$ when $x_1 < x_2$.
\begin{itemize}
\item We say a function $f$ is \underline{\textbf{strictly} increasing on the interval $I$} if for all $x_1, x_2$ in the interval, $f(x_1) < f(x_2)$ when $x_1 < x_2$.
\end{itemize}
\item We say a function $f$ is \underline{decreasing on the interval $I$} if for all $x_1, x_2$ in the interval, $f(x_1) \geq f(x_2)$ when $x_1 < x_2$.
\begin{itemize}
\item We say a function $f$ is \underline{\textbf{strictly} decreasing on the interval $I$} if for all $x_1, x_2$ in the interval, $f(x_1) > f(x_2)$ when $x_1 < x_2$.
\end{itemize}
\end{itemize}
\end{definition*}
%Definitions adapted from OpenStax Calculus Volume 1.
\vspace{0.15in}

\textbf{Exercises}: Practice with these definitions.
\begin{itemize}
\item In 3-5 sentences, describe a decreasing function (give a specific example, tell a story).
\item In 3-5 sentences, describe a decreasing function (give a specific example, tell a story).
\item In 3-5 sentences, describe a decreasing function (give a specific example, tell a story).
\item Draw a picture of a constant function.
\item Draw a picture of an increasing function.
\item Draw a picture of an increasing function that is not strictly increasing.
\item Draw a picture of a decreasing function.
\item Draw a picture of a decreasing function that is not strictly decreasing.
\item Write the equation of a constant function.
\item Write the equation of an increasing function.
\item Write the equation of a decreasing function.
\end{itemize}


 
\end{document}